All appearances whatsoever are accordingly continuous magnitudes,
either in their intuition, as extensive magnitudes, or in their mere per-
ception (sensation and thus reality), as intensive ones.

Now if all appearances, considered extensively as well as intensively, are
continuous magnitudes, then the proposition that all alteration (transi-
tion of a thing from one state into another) is also continuous could be
proved here easily and with mathematical self-evidence.

If all reality in perception has a degree, between which and negation
there is an infinite gradation of ever lesser degrees, and if likewise every
sense must have a determinate degree of receptivity for the sensations,
then no perception, hence also no experience, is possible that, whether
immediately or mediately (through whatever detour in inference one
might want), would prove an entire absence of everything real in ap-
pearance, i.e., a proof of empty space or of empty time can never be
drawn from experience.

Thus every sensation, thus also every reality in appearance, however
small it may be, has a degree, i.e., an intensive magnitude, which can
still always be diminished, and between reality and negation there is a
continuous nexus of possible realities, and of possible smaller percep-
tions.


